Liver fibrosis caused by \emph{Schistosoma japonicum} is a persistent concern in disease management in endemic regions and can be accompanied by portal hypertension and severe hepatosplenic complications \citep{colley2014schistosomiasis,mcmanus2018schistosomiasis}. In endemic settings, including China, assessment of established hepatic morbidity complements infection control \citep{mcmanus2010china,lo2022who}. Longitudinal evidence shows that severe parenchymal fibrosis can persist for years despite antiparasitic treatment \citep{carlton2010morbidity}. Repeatable imaging assessment of disease severity therefore provides an important foundation for population screening, assessment of morbidity, and long-term follow-up.

B-mode ultrasound non-invasively depicts hepatic parenchymal and portal abnormalities, providing an imaging tool suitable for repeated assessment. However, the ultrasound appearances of \emph{S. japonicum}-associated fibrosis have complex spatial patterns: interseptal fibrosis can produce network-like or fish-scale-like echoes, whereas portal abnormalities have a different anatomical distribution. The Basel ultrasonography protocol emphasizes distinguishing these patterns \citep{richter2025basel}. These features make local echoes, structural distribution, and anatomical context relevant to severity assessment. For automated grading, this raises two linked requirements: extracting information that supports fine-grained scoring from complex images, and retaining relevant spatial cues as outputs that can be compared with the original image.

Computational approaches have progressed from ultrasound texture analysis to radiomics and deep convolutional models, providing a technical foundation for reducing variation in manual interpretation \citep{yeh2003fibrosis,guo2024radiomics,lee2020ultrasound}. SFibAI specifically represents ultrasound grading of \emph{S. japonicum}-associated liver fibrosis using 36 ordered score levels from 0.0 to 3.5 at intervals of 0.1, while retaining their correspondence to the four clinical grades \citep{xu2026sfibai}. This representation provides a more detailed description of severity than coarse categories and establishes a basis for learning fine-grained scores directly from images. We retain this task representation and grading loss, use SFibAI as the direct baseline, and investigate how anatomical views and local regions can be organized and used within the scoring process.

An accurate score derived from whole-image features does not directly identify the acquisition view recognized by the model, the regions exhibiting relevant appearances, or how those regions contribute to the final prediction. Global representations can contain this information, but a score alone leaves users unable to inspect view predictions and regional responses separately. Acquisition-view labels and local lesion annotations provide complementary supervision: the former describes anatomical context, whereas the latter identifies locations of interest. Using these signals for grading requires an explicit computational connection between anatomical context, local features, and image-level judgment, while accommodating annotations that cover only part of an abnormal region.

Multitask learning can jointly learn grading and spatial tasks through shared representations \citep{caruana1997multitask}, but representation sharing alone does not specify how predicted views and localization outputs enter the scoring computation. Post-hoc attribution methods such as Grad-CAM can highlight regions associated with a prediction \citep{selvaraju2017gradcam}, while generating these visualizations leaves the original scoring pathway unchanged. We therefore ask: \emph{can acquisition views and local positions be organized as regional evidence that directly participates in prediction, improving fine-grained grading while retaining spatial outputs for human inspection?} Our design principle is to introduce view conditioning before spatial aggregation, allowing local features to form grading evidence within their anatomical context before combining them with whole-image judgment.

To implement this design, we propose VCRE-Fib (View-Conditioned Regional Evidence for Fibrosis Grading), illustrated in Figure~\ref{fig:architecture}. The framework preserves a global scoring pathway and uses view labels and local lesion boxes to learn view prediction and weak localization. It generates view-conditioned regional grading evidence before spatial pooling, uses weak localization to guide regional aggregation, adds the resulting correction to the global prediction, and mixes the conditional outputs according to predicted view probabilities. This structure explicitly connects the learning, use, and presentation of spatial information: view predictions determine the mixture of conditional outputs, localization responses participate in regional aggregation, and both remain available as inspectable outputs. Inference requires only an image and returns a fibrosis score, an acquisition view, and a localization map.

We compare VCRE-Fib with SFibAI and two independently trained module-deletion variants under a common data split and evaluation protocol. On a test set of 4,107 images from 240 patients across four centers, the full model achieves the lowest prespecified composite risk in this comparison ($\risk=0.186865$), a 7.115\% reduction relative to SFibAI, together with the lowest image-level, patient-max, patient-median, and center-balanced COR. The module-deletion comparison also shows that lower grading risk does not coincide with uniformly better spatial-task metrics, motivating joint assessment of scoring performance and spatial output quality. These results support further investigation of view-conditioned regional evidence in this task and provide a concrete framework linking fine-grained grading with inspectable spatial information. Its practical benefit for clinician review and clinical decisions remains to be evaluated.
